\documentclass{beamer}

\usepackage[T1]{fontenc}
\usepackage{graphicx}
\usepackage{hyperref}

\usetheme{Madrid}

\usepackage{helvet}
\renewcommand{\familydefault}{\sfdefault}

% --- Title Information ---
\title[Fibers of Light]{Fibers of Light: Optical Interconnects for Energy-Efficient Computing \\
\normalsize Ultra-low-power photonics, with a path toward topological architectures}
\author{Pavel Chekmaryov}
\date{\today}

\begin{document}\footnotesize


% --- Title Slide ---
\begin{frame}
    \titlepage
    \begin{center}
        \vspace{1cm}
    \large Seed Overview \\
    \normalsize \href{mailto:p@occurai.com}{p@occurai.com} \\
    \normalsize \url{https://www.linkedin.com/in/ufijuice/}
    \newline

    \end{center}
\end{frame}

% --- Slide 2: The Problem ---
\begin{frame}
    \frametitle{The Problem: Scaling Limits of Conventional Computing}
    \begin{itemize}
        \item \textbf{Bandwidth bottlenecks:} Data movement increasingly dominates system-level performance.
        \item \textbf{Energy constraints:} Performance scaling is limited by power consumption, with a large fraction of energy spent on interconnects.
        \footnote{Estimates suggest up to $\sim$40\% of datacenter power is spent on data movement.}
    \end{itemize}
\end{frame}

% --- Slide 3: The Solution ---
\begin{frame}
    \frametitle{The Solution: Photonic Interconnects as a Computing Substrate}

    \begin{block}{Our Approach}
        By leveraging nonlinear optical effects, we investigate regimes where localized optical modes can be used for structured signal propagation and processing.
    \end{block}

    \begin{block}{Core Technologies}
        \begin{itemize}
            \item \textbf{Near-term:} Soliton-based photonic interconnects for high-bandwidth, energy-efficient data movement.
            \item \textbf{Long-term:} Novel topological photonic states (e.g.\ Hopfions) as a foundation for robust, physics-enabled computation.
        \end{itemize}
    \end{block}
\end{frame}

% --- Slide 4: How It Works ---
\begin{frame}[fragile]
    \frametitle{How It Works: Physical Function and Implementation}

    \textbf{1. Nonlinear Optical Function}
        Our approach relies on controlled nonlinear interactions between optical signals.
    \begin{center}
    {\tiny
    \begin{verbatim}
      Laser A --\
                 >---[Coupler]-->[Nonlinear Region]--> Output
      Laser B --/         ^
                          |
                    [Phase Control]
    \end{verbatim}
    }
    \end{center}

    \vspace{0.4cm}

    \textbf{2. Photonic Chip Architecture}
        These nonlinear functions are embedded in a scalable silicon photonics platform
        based on networks of micro-ring resonators.
    \begin{center}
    {\tiny
    \begin{verbatim}
                                 +---[Heater]---+
                                 |       / \    |
  Input --[Waveguide]-->[Coupler]--+----O   O---+--[Coupler]--> Output
                                 |       \ /    |
                                 +---[Heater]---+
                                       ...
    \end{verbatim}
    }
    \end{center}

    \small
    \textit{Configuration elements (e.g.\ heaters) are used for tuning and routing, not computation.}
\end{frame}


% --- Slide 5: Application 1 (Solitons) ---
\begin{frame}
    \frametitle{Application 1 (Near-Term): Soliton-Based Interconnects}

    \begin{itemize}
        \item \textbf{Physical basis:} Optical solitons are stable pulses arising from a balance between dispersion and nonlinearity, allowing information to propagate with reduced distortion over distance.
        
        \item \textbf{Proposed product:} An on-chip and chip-to-chip photonic interconnect optimized for low-latency, high-bandwidth communication between processor cores and accelerators.
        
        \item \textbf{Expected impact:} By reducing signal distortion and repeater overhead, soliton-based links target lower energy per bit, reduced latency, and higher bandwidth density compared to conventional electrical interconnects.
    \end{itemize}
\end{frame}
% --- Slide 6: Topological Photonic Computing ---
\begin{frame}
    \frametitle{Exploratory R\&D: Topological Photonic Computing}
    \begin{itemize}
        \item \textbf{The Concept:} A \alert{Hopfion} is a stable, topological structure of light. Information can be stored in its robust, complex state.
        \item \textbf{The Platform:} A research setup to explore non-binary photonic computation and develop an analog module for specialized calculations.
        \item \textbf{The Potential:} Moves beyond binary systems to multi-level, analog computation. Useful for modeling complex real-world phenomena and supports \alert{Reservoir Computing}.
    \end{itemize}
\end{frame}

% --- Slide 7: Strategic Advantage ---
\begin{frame}
    \frametitle{Our Strategic Advantage}
    \begin{itemize}
        \item \textbf{Technology Readiness:} Starting with mature Silicon Nitride (SiN) prototypes, with a clear path to high-performance AlGaAs devices.
        \item \textbf{Validated Physics:} Simulations confirm the underlying models, providing a solid foundation for hardware. \textit{(Simulations are not the same as fabrication.)}
        \item \textbf{Phased Roadmap:} Step-by-step approach delivers near-term results while progressing toward new computing paradigms.
        \item \textbf{Platform Versatility:} Supports multiple computational paradigms, including Hopfions and Reservoir Computing, on the same hardware.
    \end{itemize}
\end{frame}

% --- Slide 8: Phased Roadmap ---
\begin{frame}[allowframebreaks]
    \frametitle{Phased Roadmap \& Integrated R\&D}

    \begin{block}{Phase 1: Soliton Accelerator (0-2 Years)}
        \begin{itemize}
            \item Deliver a working \textbf{Soliton Accelerator} for inter-core communication.
            \item Publish performance benchmarks compared to electrical interconnects.
            \item \textit{R\&D Focus: Begin integration of \textbf{on-chip light sources} (optional but high-value).}
        \end{itemize}
    \end{block}

    \begin{block}{Phase 2: Hopfion Compute Module (2-5 Years)}
        \textit{Conditional on successful completion of Phase 1.}
        \begin{itemize}
            \item Demonstrate stable creation and control of single \textbf{Hopfions}.
            \item Develop the first integrated \textbf{Hopfion Compute Module}.
        \end{itemize}
    \end{block}

    \framebreak

    \begin{block}{Phase 3: Photonic Motherboard (5-8 Years)}
        \textit{Long-term scientific vision (not part of current raise).}
        \begin{itemize}
            \item Expand the photonic fabric to replace major inter-chip communication (CPU-GPU-RAM).
            \item \textit{R\&D Focus: Mature \textbf{on-chip light sources} for multi-chip integration.}
        \end{itemize}
    \end{block}

    \begin{alertblock}{Ultimate Vision: Fibration Computer (8+ Years)}
        Merge all computation into a unified \textbf{Hopf fibration structure}, where microresonators act as the "fibers" of the system.
    \end{alertblock}

\end{frame}

% --- Slide 9: The Ask ---
\begin{frame}
    \frametitle{The Ask}
    We are raising a \textbf{$2,000,000 Seed round} to achieve key milestones over the next 18 months:
    \begin{itemize}
        \item \textbf{Team:} Hire 1-2 photonic engineers and secure key advisors.
        \item \textbf{Lab \& PoC:} Set up a lab and build the first physical Proof-of-Concept device to validate simulations.
        \item \textbf{Fabrication:} Complete at least two design-and-fabrication cycles for the Soliton Accelerator chip.
        \item \textbf{Runway:} Provide ~18 months of runway to demonstrate hardware feasibility and prepare for a larger seed round.
    \end{itemize}
\end{frame}

% --- Slide 10: Founder \& Advisors ---
\begin{frame}
    \frametitle{Founder \& Advisors}

    \begin{block}{Pavel Chekmaryov — Founder \& CEO}
        \begin{itemize}
            \item Experience in both hardware and software development.
            \item Worked at \textbf{Snowcap Compute}, a quantum hardware startup.
            \item Background in technical recruitment and early-stage startup operations.
            \item Hands-on experience with product development and startup growth.
        \end{itemize}
    \end{block}

    \begin{block}{Scientific Advisors}
        \begin{itemize}
            \item \textbf{Anton Velichko (University of Sheffield)} — semiconductor optics and spectroscopy, with focus on material realism and experimental feasibility.
            \item Expertise in \textbf{AlGaAs} and semiconductor physics.
            \item Guidance on fabrication constraints and theory–experiment alignment.
        \end{itemize}
    \end{block}

    \begin{block}{Team Expansion}
        \small\textit{Exploring collaborations with technical co-founders and domain experts as the project develops.}
    \end{block}

\end{frame}

% --- Slide 11: Competitive Landscape ---
\begin{frame}
    \frametitle{Competitive Landscape}
    Our strategy bridges near-term interconnect improvements and long-term novel computing.

    \centering
    \resizebox{0.9\textwidth}{!}{
    \begin{tabular}{|l|p{5cm}|p{5cm}|}
        \hline
        \textbf{Category} & \textbf{Existing Players / Paradigm} & \textbf{New Approach / Vision} \\
        \hline
        Core Computing & \begin{itemize}
            \item Nvidia (GPUs)
            \item Google (TPUs)
        \end{itemize} & \begin{itemize}
            \item \textbf{Hopfion Architecture} (Analog, early-stage)
            \item PsiQuantum, Xanadu (Quantum)
        \end{itemize} \\
        \hline
        Interconnects & \begin{itemize}
            \item Electrical Interconnects
            \item Ayar Labs, Lightmatter (optical interconnect competitors)
        \end{itemize} & \begin{itemize}
            \item \textbf{Soliton Accelerator} (our technology)
        \end{itemize} \\
        \hline
    \end{tabular}
    }

    \vspace{0.5cm}
    \begin{alertblock}{}
        We begin by advancing interconnects with our \textbf{Soliton Accelerator}, then use this platform to enable a new computing paradigm with the \textbf{Hopfion Architecture}.
    \end{alertblock}
\end{frame}

% --- Slide 12: Market Opportunity ---
\begin{frame}
    \frametitle{Market Opportunity}
    Addressing bottlenecks in two fast-growing technology markets.

    \begin{block}{Target Market 1: Data Center \& HPC Interconnects}
        \begin{itemize}
            \item Optical interconnect market projected to exceed \textbf{\$27B by 2027}.
            \item Data movement consumes a large portion of data center energy—our \textbf{Soliton Accelerator} addresses this.
            \item \textbf{Initial Customers:} Hyperscalers with high-performance, energy-efficient interconnect needs.
        \end{itemize}
    \end{block}

    \begin{block}{Target Market 2: Next-Generation Computing Hardware}
        \begin{itemize}
            \item AI hardware market projected at \textbf{\$150B+ by 2027}.
            \item Positions us to capture part of the post-GPU computing paradigm.
        \end{itemize}
    \end{block}
\end{frame}

% --- Final Slide ---
\begin{frame}
    \frametitle{Thank You}
    \begin{center}
        \Huge Thank You
        \vspace{1cm}

        \Large Pavel Chekmaryov \\
        \normalsize \href{mailto:p@occurai.com}{p@occurai.com} \\
        \normalsize \url{https://www.linkedin.com/in/ufijuice/}
    \end{center}
\end{frame}

\end{document}